%% kapa: exam=1 round=21 year=2010 period=1 form=A subject=civil_law subject_ko=민법 seq=29 number=29
%% source: Q-Net 감정평가사 기출문제 (cst00309) artlSeq=5024181 file=extracted/1008/5024181/2010년 제21회 감정평가사 1차시험 기출문제/1교시.hwp
%% conversion: kordoc 4.24.51 -> tools/gen1.py (auto); review=proofread-sonnet; flags=none
\begin{question}{29}
유치권에 관한 설명으로 옳지 않은 것은? (다툼이 있으면 판례에 의함)
\begin{choices}
  \item 당사자 사이에 유치권의 성립을 배제하기로 하는 약정은 유효하다.
  \item 임차인의 부속물매수청구권에 따른 대금채권을 담보하기 위한 임대목적물에 대한 유치권은 인정되지 않는다.
  \item 유치물의 멸실, 훼손으로 인하여 소유자가 받을 금전 그 밖의 물건에 대하여는 유치권의 효력이 미치지 않는다.
  \item 임차인은 권리금반환청구권을 담보하기 위해 임차목적물에 유치권을 행사할 수 없지만, 보증금반환청구권을 담보하기 위해서는 임차목적물에 유치권을 행사할 수 있다.
  \item 유치권자가 유치물에 대한 보존행위로서 목적물을 사용한 경우, 불법행위로 인한 손해배상책임을 지지 않는 것이 원칙이다.
\end{choices}
\end{question}
\answer{4}
